% SPDX-License-Identifier: CC-BY-SA-4.0
% Author: Matthieu Perrin
% Part: <Nom de la partie>
% Section: <Nom de la section>
% Sub-section: <Nom de la sous-section>  % (facultatif, laisser vide si non utilisé)
% Frame: <Titre de la slide>

\begingroup

\begin{frame}{Ambiguïté des grammaires algébriques}
  
  \Probleme{Ambiguity}{
    Une grammaire algébrique $G$. 
  }{
    Y a-t-il un mot $w$ avec deux arbres de dérivation ?
  }

  \begin{exampleblock}{Exemples -- les grammaires suivantes sont-elles ambiguës ?}

    \begin{itemize}
    \item $G_1 = \langle \{a, +\}, \{A\}, A, \left\{\example{A \rightarrow A + a \mid a}\right\} \rangle \uncover<2->{\alert{\quad\in \textsc{neg}(\textsc{Ambiguity})}}$
    \item $G_2 = \langle \{a, +\}, \{A\}, A, \left\{\example{A \rightarrow A + A \mid a}\right\} \rangle \uncover<2->{\example{\quad\in \textsc{pos}(\textsc{Ambiguity})}} $
    \end{itemize}

    \vspace{2mm}
    \uncover<2->{
      Le mot $\example{a+a+a}$ a deux arbres de dérivation selon $G_2$ :

      ~\hfill
      \begin{tikzpicture}[tree, x=8mm, y=6.5mm, tree node internal/.append style={alert,}, tree node leave/.append style={example,},]
        \tree{$A$}{
          \tree{$A$}{
            \tree{$A$}{\tree{$a$}{}}
            \tree{$+$}{}
            \tree{$A$}{\tree{$a$}{}}
          }
          \tree{$+$}{}
          \tree{$A$}{\tree{$a$}{}}
        }
      \end{tikzpicture}
      \hfill
      \begin{tikzpicture}[tree, x=8mm, y=6.5mm, tree node internal/.append style={alert,}, tree node leave/.append style={example,},]
        \tree{$A$}{
          \tree{$A$}{\tree{$a$}{}}
          \tree{$+$}{}
          \tree{$A$}{
            \tree{$A$}{\tree{$a$}{}}
            \tree{$+$}{}
            \tree{$A$}{\tree{$a$}{}}
          }
        }
      \end{tikzpicture}
      \hfill~
    }
  \end{exampleblock}    

\end{frame}

\endgroup

